\section{Details of the cluster computation}\label{app:cluster}

The CPU partition of YUCA comprises $24$ nodes with two $32$-core AMD EPYC 9354 processors per node, $1$\,TB of RAM and no hyper-threading. The archived system records describe Rocky Linux~9.6, SLURM~24.11 and GCC~11. The order-$13$ run used \texttt{nauty}~2.7r4; the valid $K_4$-free order-$14$ run used the locally compiled version~2.8. The two projects each had a $500$-core quality-of-service limit. Algorithm~\ref{alg:dist} summarises the scheduling design; the qualifications below describe the actual concurrent workers.

\begin{algorithm}[t]
\caption{Distributed search on the cluster}\label{alg:dist}
\KwIn{$n$; number of slices $M$; number of array tasks $T$; \texttt{geng} options $\Omega$; mode (stuck / table)}
A one-core wrapper job compiles the program, creates the registries \texttt{listas\_t}, \texttt{enc\_t} ($t=0,\dots,T-1$) and submits a SLURM array of $T$ tasks with $64$ cores each\;
\ForEach{task $t$ in parallel}{
  \ForEach{slice $r\equiv t\pmod T$, $0\le r<M$, with $64$ slices in flight (\texttt{xargs -P 64})}{
    \lIf{$r\in$ \texttt{listas\_t}}{skip (resumable)}
    run \texttt{geng $\Omega$ $n$ -q $r/M$ | atascadas -{}-modo} mode, writing the output to a per-slice file\;
    under the lock used by that worker, append the output to \texttt{enc\_t} and $r$ to \texttt{listas\_t}\;
  }
}
The wrapper polls progress every $60$\,s; when all tasks end it sums the counters, verifies that $\bigcup_t\texttt{listas\_t}=\{0,\dots,M-1\}$, lists duplicated slices, and records \texttt{sacct} accounting\;
\end{algorithm}

\paragraph{Coverage and counters}
The seven raw completion registries for each run were checked against their prescribed ranges and congruence classes. The order-$12$ run contains all $\num{4000}$ indices exactly once, and the order-$13$ run all $\num{200000}$ indices exactly once. For order $14$, the $\num{2000692}$ entries cover all $\num{2000000}$ indices, with $692$ occurring twice and none more often. The distinct-index file and duplicated-index list in \protect\path{yuca/logs/} agree with this recomputation. Per-directory summaries contain all $200$ directories at order $13$ and all $\num{2000}$ at order $14$, with $\num{1000}$ result files each. Summing their counters once per slice reproduces Table~\ref{tab:runs}. After the edge-neighbourhood filter, $\num{218516}$ graphs at order $13$ and $\num{13536597}$ at order $14$ reach the full homology test; the respective outputs contain one and $51$ graphs.

\paragraph{Accounting}
The valid order-$12$ array is 2079051 (wrapper 2079049), and the order-$13$ array is 2079059 (wrapper 2079058). Their task spans are $7$ min $10$ s and $7$ h $10$ min $11$ s. The valid order-$14$ arrays, 2079102 on p029 and 2079156 on p011, span 11 October 2026, 11:23:46, to 12 October, 22:26:34, in Hermosillo time (UTC$-7$): $35$ h $2$ min $48$ s. Up to $896$ cores were allocated, with the second front starting later. Allocated core time must not be interpreted as measured CPU time or as a separate benchmark for the generator and filter.

The cardinality rows were computed by separate runs in table mode. The second order-$12$ run (array 2084878, $\num{4000}$ slices of \texttt{geng -c 12}, $448$ cores, $8$ min) computed all columns and reproduced the six columns of the first run exactly. The order-$13$ row used arrays 2084597 (p029) and 2084636 (p011), with $\num{200000}$ slices of \texttt{geng -c 13} on up to $896$ cores, from 12 October, 21:04, to 13 October, 17:59 (UTC$-7$). Its registries cover all slices, $483$ of them twice. The connected counters, summed once per slice, equal the number of connected graphs on $13$ vertices (OEIS A001349). This agreement is a consistency check. It would show that no slice is missing or truncated only under premises that the archived run does not certify: that the residue classes partition the connected graphs, that the stream of each slice contains each graph of its class at most once and no other graph, that the counters are intact, and that exactly one result is summed per residue. Under these premises, writing $N_r$ for the size of class $r$ and $n_r$ for the count observed, $0\le n_r\le N_r$ and $\sum_r n_r=\sum_r N_r$ force $n_r=N_r$ for every $r$; without them, omissions and duplications could compensate each other in the total. The recomputation described below checks these premises slice by slice.

\paragraph{Correction to the archived worker}
In commit \texttt{4a059b5} the file named \protect\path{yuca/as_run/tarea_n14_p011.sh} has SHA-256
\begin{center}\scriptsize\texttt{50477550293c8442ee6368b360dbeade20144fe8f20be54f32571d263018b320}\end{center}
and belongs to the earlier failed attempt with \texttt{geng}~2.7r4 (array 2079068), not the successful descending p011 front. Since commit \texttt{892e41b} that file is kept unchanged as \protect\path{yuca/as_run/tarea_n14_p029_intento_fallido_2079068.sh}, and the name \protect\path{tarea_n14_p011.sh} holds the worker identified for the successful front, found in the primary provenance archive; its $927$ bytes have SHA-256
\begin{center}\scriptsize\texttt{b95a125b466b47e921269baa43028983ea00081563ae3d4a126fc86452381ba9}.\end{center}
A copy accompanies the manuscript sources as \protect\path{appendix_data/tarea_n14_p011_correct.sh}. SLURM accounting on YUCA does not retain submitted batch scripts, so the identification is a forensic corroboration from file-system records, not the recovery of an authenticated copy of the submitted script: this file was created at the second at which array 2079156 was submitted (11 October 2026, 15:40:09, UTC$-7$) and has not been modified since (equal birth, modification and change times); the task logs of that array print the descending starting residues that this script computes; and no file in the directory of the valid run predates that run, while the failed attempt is preserved in a separate directory. The failed run cannot be certified by its SLURM completion status alone: it lacked pipeline error propagation for the unsupported generator option.

\paragraph{Duplicated slices and remaining certification work}
The two valid fronts shared completion registries and per-slice result paths, but used different lock files. Consequently the recorded coverage is complete while the archive does not preserve two independently comparable outputs for each duplicate. The assertion that all duplicated outputs were identical is therefore not retained. Two recomputations close these gaps: a complete rerun of the order-$14$ search with a pipeline that checks the exit status and the graph count of both processes of every slice and publishes one receipt per slice atomically, and complete-state recomputations of the order-$12$ and order-$13$ rows. Until they are completed, the completeness claims remain conditional on the stated hypotheses and are not presented as fully certified computations. No error in the displayed finite certificates has been inferred from these implementation limitations.
